% --- section_1.tex ---
\begin{tabularx}{\columnwidth}{XX}
  \multicolumn{2}{l}{\textbf{基本概念}} \\
  瞬时速度\newline {\footnotesize 写出极限定义式（位移对时间的极限）。} & \\
  平均速度\newline {\footnotesize 写出位移与所用时间的比值（矢量）。} & \\
  瞬时速率\newline {\footnotesize 写出瞬时速度的大小（标量）。} & \\
  平均速率\newline {\footnotesize 写出路程与所用时间的比值。} & \\
  平均加速度\newline {\footnotesize 写出速度变化量与时间的比值。} & \\
  \multicolumn{2}{l}{\textbf{匀变速直线运动}} \\
  位移\newline {\footnotesize 匀变速直线运动：位移与时间、初速度、加速度的关系式。} & \\
  {\footnotesize 同一过程：用平均速度表示的位移公式。} & \\
  速度\newline {\footnotesize 匀变速直线运动：速度与时间、初速度、加速度的关系式。} & \\
  速度平方\newline {\footnotesize 匀变速直线运动：速度与位移关系（不含时间）。} & \\
  中间位移速度\newline {\footnotesize 匀变速过程：位移中点处的瞬时速度公式。} & \\
  中间时刻速度\newline {\footnotesize 匀变速过程：时间中点对应的瞬时速度（平均速度）。} & \\
  相邻相等时间间隔位移差\newline {\footnotesize 匀变速：连续相等时间 $T$ 内位移之差与加速度关系。} & \\
  第M段与第N段Ts时间内位移差\newline {\footnotesize 匀变速：第 $M$ 段与第 $N$ 段（每段时长 $T$）位移之差。} & \\
  \multicolumn{2}{l}{\textbf{从静止开始加速}} \\
  相等时间每段位移比\newline {\footnotesize 从静止匀加速：各 $T$ 内位移之比。} & \\
  相等时间总计位移比\newline {\footnotesize 从静止匀加速：前 $n$ 个 $T$ 总位移之比。} & \\
  相等位移间隔每段用时比\newline {\footnotesize 从静止匀加速：各相等位移段所用时间之比。} & \\
  相等位移间隔总计用时比\newline {\footnotesize 从静止匀加速：通过前 $n$ 段相等位移累计用时之比。} & \\
\end{tabularx}
\medskip

% --- section_2.tex ---
\begin{tabularx}{\columnwidth}{XX}
  水平位移 & \\
  竖直位移 & \\
  竖直速度\newline {\footnotesize 平抛：竖直分速度与下落时间或竖直位移的关系。} & \\
  合速度\newline {\footnotesize 平抛：合速度大小（水平、竖直分量合成）。} & \\
  轨迹抛物线方程\newline {\footnotesize 消去时间：$y$ 与 $x$ 的抛物线关系。} & \\
  \multicolumn{2}{l}{\textbf{平抛运动推论}} \\
  速度偏角\newline {\footnotesize 速度与水平方向夹角的正切（可用 $v_y/v_0$、时间或位移表示）。} & \\
  斜面抛回斜面\newline {\footnotesize 落到斜面时：位移偏角等于斜面倾角，写出由此得到的关系。} & \\
  垂直打到斜面上\newline {\footnotesize 垂直撞击斜面时速度偏角与斜面倾角的关系式。} & \\
\end{tabularx}
\medskip

% --- section_3.tex ---
\begin{tabularx}{\columnwidth}{XX}
  向心力\newline {\footnotesize 用线速度写向心力大小（圆周运动）。} & \\
  向心力\newline {\footnotesize 用向心加速度写向心力（$F=ma_{\text{向}}$）。} & \\
  线速度\newline {\footnotesize 圆周运动：线速度与角速度、半径关系。} & \\
  角速度\newline {\footnotesize 圆周运动：角速度与周期关系。} & \\
\end{tabularx}
\medskip

% --- section_4.tex ---
\begin{tabularx}{\columnwidth}{XX}
  球体万有引力\newline {\footnotesize 均匀球体：体外与体内引力大小的分段表达式（含条件）。} & \\
  球体万有引力势能\newline {\footnotesize 质点在球外时的引力势能（规定无穷远为零势能）。} & \\
  向心力简化式\newline {\footnotesize 圆轨道：由线速度写 $GM=v^2 r$。} & \\
  {\footnotesize 圆轨道：由角速度写 $GM=\omega^2 r^3$。} & \\
  {\footnotesize 圆轨道：由向心加速度写 $GM=a r^2$。} & \\
  {\footnotesize 圆轨道：由周期写 $GM=4\pi^2 r^3/T^2$。} & \\
  密度公式\newline {\footnotesize 近地卫星测密度：由周期写平均密度。} & \\
  \multicolumn{2}{l}{\textbf{开普勒三大定律}} \\
  周期定律\newline {\footnotesize 开普勒第三定律：半长轴立方与周期平方之比与中心天体质量。} & \\
  \multicolumn{2}{l}{\textbf{黄金代换}} \\
  黄金代换\newline {\footnotesize 地面附近：引力与重力近似时的 $GM$ 与 $g、R$ 关系。} & \\
  \multicolumn{2}{l}{\textbf{宇宙速度}} \\
  第一宇宙速度\newline {\footnotesize 贴近星球表面圆轨道环绕速度。} & \\
  \multicolumn{2}{l}{\textbf{追及问题}} \\
  同向（最近到最远）\newline {\footnotesize 同向卫星：从最近到最远经历半圈相对转角所需时间。} & \\
  同向（最近到最近）\newline {\footnotesize 同向卫星：从最近到再次最近相对转角 $2\pi$ 所需时间。} & \\
  反向（最近到最远）\newline {\footnotesize 反向卫星：从最近到最远相对转角 $\pi$ 所需时间。} & \\
  反向（最近到最近）\newline {\footnotesize 反向卫星：从最近到再次最近相对转角 $2\pi$ 所需时间。} & \\
  \multicolumn{2}{l}{\textbf{双星问题}} \\
  质量半径关系\newline {\footnotesize 双星：两星到质心的距离与质量成反比。} & \\
  质量角速度关系\newline {\footnotesize 双星：角速度与间距、总质量的关系式。} & \\
\end{tabularx}
\medskip

% --- section_5.tex ---
\begin{tabularx}{\columnwidth}{XX}
  重力\newline {\footnotesize 地表附近重力大小（质量与 $g$）。} & \\
  两极重力加速度\newline {\footnotesize 自转影响：极点视重与万有引力关系（无向心加速度）。} & \\
  赤道重力加速度\newline {\footnotesize 自转影响：赤道视重（万有引力减自转向心力）。} & \\
  滑动摩擦力\newline {\footnotesize 写滑动摩擦定律（正压力与动摩擦因数）。} & \\
  最大静摩擦力\newline {\footnotesize 写出最大静摩擦与正压力关系（静摩擦因数）。} & \\
  弹簧弹力\newline {\footnotesize 胡克定律（伸长或压缩时的弹力）。} & \\
  共点力平衡\newline {\footnotesize 共点力平衡的矢量条件（合力为零）。} & \\
\end{tabularx}
\medskip

% --- section_6.tex ---
\begin{tabularx}{\columnwidth}{XX}
  \multicolumn{2}{l}{\textbf{牛顿运动定律}} \\
  牛顿第二定律\newline {\footnotesize 合外力与加速度的关系（矢量式或分量式）。} & \\
  牛顿第三定律\newline {\footnotesize 作用力与反作用力：等大、反向、共线，作用在两个物体上。} & \\
  重力\newline {\footnotesize 动力学中重力大小（与静力学表述一致即可）。} & \\
  滑动摩擦力\newline {\footnotesize 写滑动摩擦定律（正压力与动摩擦因数）。} & \\
  最大静摩擦力\newline {\footnotesize 写出最大静摩擦与正压力关系（静摩擦因数）。} & \\
  胡克定律\newline {\footnotesize 弹簧弹力与形变量（劲度系数）关系（注明伸长或压缩）。} & \\
  \multicolumn{2}{l}{\textbf{空气阻力}} \\
  空气阻力\newline {\footnotesize 题给阻力模型（如与速度成正比或平方）时的表达式。} & \\
\end{tabularx}
\medskip

% --- section_7.tex ---
\begin{tabularx}{\columnwidth}{XX}
  动能定理\newline {\footnotesize 合外力做功等于动能增量（动能定理）。} & \\
  动能\newline {\footnotesize 动能定义式 $\frac12 mv^2$。} & \\
  重力做功\newline {\footnotesize 重力做功与初末高度差的关系（或下落高度）。} & \\
  弹力做功\newline {\footnotesize 弹簧弹力做功与形变量的关系；若表中为“弹力与速度垂直”特例请按该条件书写。} & \\
\end{tabularx}
\medskip

% --- section_8.tex ---
\begin{tabularx}{\columnwidth}{XX}
  功\newline {\footnotesize 恒力做功：力、位移及夹角余弦。} & \\
  合力做功\newline {\footnotesize 合力做功等于动能增量（动能定理）。} & \\
  动能\newline {\footnotesize 动能定义式 $\frac12 mv^2$。} & \\
  重力势能 & \\
  弹性势能 & \\
  重力做功与重力势能\newline {\footnotesize 重力做功等于重力势能增量的负值。} & \\
  弹力做功与弹性势能\newline {\footnotesize 弹力做功等于弹性势能增量的负值。} & \\
  机械能 & \\
  机械能守恒 & \\
  非保守力做功\newline {\footnotesize 功能关系：非保守力做功与机械能变化。} & \\
\end{tabularx}
\medskip

% --- section_9.tex ---
\begin{tabularx}{\columnwidth}{XX}
  动量定理\newline {\footnotesize 合外力冲量等于动量变化（矢量式）。} & \\
  恒力冲量\newline {\footnotesize 冲量：力与作用时间的乘积（矢量）。} & \\
  平均力冲量\newline {\footnotesize 用平均力表示的冲量（动量变化）。} & \\
  系统动量定理\newline {\footnotesize 系统所受合外力的冲量等于系统总动量变化。} & \\
\end{tabularx}
\medskip

% --- section_30.tex ---
\begin{tabularx}{\columnwidth}{XX}
  动量守恒\newline {\footnotesize 系统所受合外力为零时总动量守恒的表达式。} & \\
  恢复系数\newline {\footnotesize 碰撞恢复系数定义（分离与接近的相对速度大小之比）。} & \\
  一维恢复系数\newline {\footnotesize 对心碰撞：分离相对速度与接近相对速度之比。} & \\
  约化质量\newline {\footnotesize 两体相对运动：约化质量 $\mu$ 与两质量关系。} & \\
  碰撞损失机械能\newline {\footnotesize 完全非弹性或一般碰撞：损失的动能表达式。} & \\
\end{tabularx}
\medskip

% --- section_24.tex ---
\begin{tabularx}{\columnwidth}{XX}
  回复力\newline {\footnotesize 简谐振动：回复力与位移的比例关系（指向平衡位置）。} & \\
  简谐振动方程\newline {\footnotesize 位移随时间的余弦（或正弦）形式。} & \\
  速度与加速度\newline {\footnotesize 简谐振动：速度、加速度与位移或相位的关系。} & \\
  周期频率关系\newline {\footnotesize 周期、频率与角频率之间的关系。} & \\
  弹簧振子周期\newline {\footnotesize 理想弹簧振子振动周期公式。} & \\
  单摆回复力\newline {\footnotesize 小角度下单摆沿弧切向的回复力与角位移关系。} & \\
  单摆周期\newline {\footnotesize 小角度单摆周期公式。} & \\
  水平弹簧振子能量\newline {\footnotesize 弹簧振子：动能与弹性势能之和（守恒式）。} & \\
\end{tabularx}
\medskip

% --- section_25.tex ---
\begin{tabularx}{\columnwidth}{XX}
  波动通式\newline {\footnotesize 沿 $x$ 轴传播的平面简谐波位移表达式。} & \\
  任意点振动方程\newline {\footnotesize 已知波动方程写某点的振动方程（令 $x$ 为常数）。} & \\
  波速波长频率关系\newline {\footnotesize $v、\lambda、f、\omega、k$ 之间的关系。} & \\
  相位差\newline {\footnotesize 两点在同一时刻的相位之差（与波程差联系）。} & \\
  \multicolumn{2}{l}{\textbf{两波源在该点起振方向相同}} \\
  干涉加强条件\newline {\footnotesize 两波源起振方向相同：振动加强的波程差条件。} & \\
  干涉减弱条件\newline {\footnotesize 两波源起振方向相同：振动减弱条件。} & \\
  \multicolumn{2}{l}{\textbf{两波源在该点起振方向相反}} \\
  干涉加强条件\newline {\footnotesize 两波源起振方向相反：振动加强的波程差条件。} & \\
  干涉减弱条件\newline {\footnotesize 两波源起振方向相反：振动减弱条件。} & \\
  衍射条件 & 障碍尺寸越小、波长越大，衍射越明显。 \\
  驻波条件 & 两列同频同振幅机械波相向传播 \\
  多普勒效应 & 相互靠近频率变高；相互远离频率变低。 \\
\end{tabularx}
\medskip

% --- section_10.tex ---
\begin{tabularx}{\columnwidth}{XX}
  库仑定律 & \\
  电场强度定义\newline {\footnotesize 试探电荷所受电场力与电荷量的比值（矢量）。} & \\
  点电荷场强\newline {\footnotesize 真空点电荷电场强度大小（库仑定律导出）。} & \\
  电场叠加\newline {\footnotesize 多个电荷产生的电场强度矢量叠加原理。} & \\
  匀强电场\newline {\footnotesize 匀强电场强度与沿场强方向电势差、距离关系。} & \\
\end{tabularx}
\medskip

% --- section_31.tex ---
\begin{tabularx}{\columnwidth}{XX}
  电势定义\newline {\footnotesize 电场中某点电势（相对零势点）的定义式。} & \\
  电势差定义\newline {\footnotesize 两点间电势差与电场力移送电荷做功关系。} & \\
  电场力做功\newline {\footnotesize 电场力做功与初末电势差、电荷量关系。} & \\
  电势能变化\newline {\footnotesize 电场力做功与电势能变化的关系。} & \\
  匀强电场电势差\newline {\footnotesize 匀强场：沿任意方向的电势差与场强投影、距离关系。} & \\
  电场力做功\newline {\footnotesize 电场力做功与初末电势差、电荷量关系。} & \\
\end{tabularx}
\medskip

% --- section_32.tex ---
\begin{tabularx}{\columnwidth}{XX}
  电容定义\newline {\footnotesize 电容器电容：电荷量与两极板电势差之比。} & \\
  平行板电容器\newline {\footnotesize 真空平行板电容与面积、间距关系。} & \\
  电容器能量\newline {\footnotesize 电容器储能公式（用电容与电压表示）。} & \\
  加速电场\newline {\footnotesize 带电粒子经加速电场获得的动能（初速为零时常用形式）。} & \\
  偏转电场加速度\newline {\footnotesize 偏转板内垂直初速方向的加速度（不计重力）。} & \\
  偏转位移\newline {\footnotesize 类平抛：穿出偏转板时在偏转方向的位移公式。} & \\
\end{tabularx}
\medskip

% --- section_11.tex ---
\begin{tabularx}{\columnwidth}{XX}
  电流定义式\newline {\footnotesize 电流：单位时间通过截面的电荷量。} & \\
  电流微观表达式\newline {\footnotesize 金属导体：电流与载流子数密度、电荷、漂移速率、截面积。} & \\
  电阻定义式\newline {\footnotesize 欧姆定律：电阻与电压、电流关系。} & \\
  电阻决定式\newline {\footnotesize 电阻定律：与电阻率、长度、截面积关系。} & \\
  串联电阻\newline {\footnotesize 串联总电阻。} & \\
  并联电阻\newline {\footnotesize 并联总电阻（两电阻倒数之和形式亦可）。} & \\
\end{tabularx}
\medskip

% --- section_33.tex ---
\begin{tabularx}{\columnwidth}{XX}
  电动势\newline {\footnotesize 电源电动势定义（非静电力移送电荷所做功与电荷量之比）。} & \\
  闭合电路欧姆定律\newline {\footnotesize 全电路电流与电动势、内外阻关系。} & \\
  路端电压\newline {\footnotesize 路端电压与电流、电动势、内阻关系。} & \\
  纯电阻功率\newline {\footnotesize 纯电阻电路功率（任选一种常用形式）。} & \\
  电源总功率\newline {\footnotesize 电源消耗化学能等的功率 $EI$。} & \\
  电源输出功率\newline {\footnotesize 电源输出到外电路的功率。} & \\
  电源效率\newline {\footnotesize 输出功率与总功率之比。} & \\
\end{tabularx}
\medskip

% --- section_34.tex ---
\begin{tabularx}{\columnwidth}{XX}
  电流表改装分流电阻\newline {\footnotesize 扩大量程：并联分流电阻与满偏电流、量程关系。} & \\
  电压表改装串联电阻\newline {\footnotesize 扩大量程：串联分压电阻关系。} & \\
  欧姆表回路\newline {\footnotesize 欧姆表：闭合回路欧姆定律形式（含调零电阻模型）。} & \\
  欧姆表中值电阻\newline {\footnotesize 指针半偏时待测电阻与中值电阻关系。} & \\
  电桥平衡条件\newline {\footnotesize 惠斯通电桥平衡时四个电阻的比例关系。} & \\
\end{tabularx}
\medskip

% --- section_12.tex ---
\begin{tabularx}{\columnwidth}{XX}
  磁感应强度\newline {\footnotesize 电流元受力定义式或运动电荷受力给出的 $B$（写出一种）。} & \\
  安培力\newline {\footnotesize 通电导线在磁场中所受安培力大小（电流与磁场垂直时）。} & \\
  通电导线受力方向 & 左手定则 \\
  直导线磁场方向 & 安培定则 \\
  螺线管磁场方向 & 右手螺旋定则 \\
\end{tabularx}
\medskip

% --- section_35.tex ---
\begin{tabularx}{\columnwidth}{XX}
  洛伦兹力\newline {\footnotesize 运动电荷在磁场中所受洛伦兹力大小（$v\perp B$）。} & \\
  向心力关系\newline {\footnotesize 带电粒子匀速圆周：洛伦兹力提供向心力。} & \\
  半径\newline {\footnotesize 匀强磁场中圆轨道半径与 $mv、qB$ 关系。} & \\
  周期\newline {\footnotesize 带电粒子在匀强磁场中圆周运动周期。} & \\
  角速度\newline {\footnotesize 圆周运动：角速度与周期关系。} & \\
  运动时间\newline {\footnotesize 转过圆心角 $\theta$ 时在磁场中运动时间与周期关系。} & \\
\end{tabularx}
\medskip

% --- section_36.tex ---
\begin{tabularx}{\columnwidth}{XX}
  速度选择器\newline {\footnotesize 电场与磁场平衡时能以直线穿过的速度条件。} & \\
  选择速度\newline {\footnotesize 速度选择器中选出的速度大小（$E$ 与 $B$ 表示）。} & \\
  加速电场\newline {\footnotesize 带电粒子经加速电场获得的动能（初速为零时常用形式）。} & \\
  质谱仪半径\newline {\footnotesize 经速度选择后粒子在磁场中的偏转半径公式。} & \\
  霍尔电压\newline {\footnotesize 霍尔效应：横向电势差与 $I、B、nq、d$ 等关系（模型依题）。} & \\
  回旋加速器最大速度\newline {\footnotesize 由磁场与 $D$ 形盒半径限制的最大回旋速度。} & \\
  回旋加速器最大动能\newline {\footnotesize 与最大速度对应的动能。} & \\
\end{tabularx}
\medskip

% --- section_13.tex ---
\begin{tabularx}{\columnwidth}{XX}
  磁通量\newline {\footnotesize 匀强磁场穿过平面的磁通量定义（面积与法线夹角）。} & \\
  法拉第电磁感应定律\newline {\footnotesize 感应电动势大小与磁通量变化率关系。} & \\
  感应电流\newline {\footnotesize 闭合回路欧姆定律：感应电流与感应电动势、电阻。} & \\
  回路电荷量\newline {\footnotesize 感应过程中通过回路截面的电荷量与磁通变化、电阻关系。} & \\
  导体棒切割\newline {\footnotesize 棒长 $l$、速度 $v$、磁场 $B$ 互相垂直时的动生电动势。} & \\
\end{tabularx}
\medskip

% --- section_37.tex ---
\begin{tabularx}{\columnwidth}{XX}
  动生电动势\newline {\footnotesize 同上（单棒模型）。} & \\
  回路电流\newline {\footnotesize 棒电阻与负载构成回路时的电流。} & \\
  安培力\newline {\footnotesize 通电导线在磁场中所受安培力大小（电流与磁场垂直时）。} & \\
  安培力冲量\newline {\footnotesize 安培力对棒的冲量与电荷量、杆长、磁场关系（常见积分结论）。} & \\
  感应电流功率\newline {\footnotesize 电路中电功率（或力的功率）与感应电流关系。} & \\
  恒力最终速度\newline {\footnotesize 棒受恒外力与安培力平衡时的收尾速度。} & \\
\end{tabularx}
\medskip

% --- section_38.tex ---
\begin{tabularx}{\columnwidth}{XX}
  磁通量\newline {\footnotesize 匀强磁场穿过平面的磁通量定义（面积与法线夹角）。} & \\
  线圈转动电动势\newline {\footnotesize 线圈在匀强磁场中转动产生的瞬时电动势（从中性面计时）。} & \\
  最大电动势\newline {\footnotesize 交流发电机峰值电动势。} & \\
  从中性面计时\newline {\footnotesize 写出从中性面开始计时的电动势瞬时值（正弦形式）。} & \\
  周期与角速度\newline {\footnotesize 交流电周期与角频率关系。} & \\
  双棒动量关系\newline {\footnotesize 光滑导轨双棒：安培力等大反向时常用的系统动量守恒式。} & \\
\end{tabularx}
\medskip

% --- section_14.tex ---
\begin{tabularx}{\columnwidth}{XX}
  双缝干涉条纹间距\newline {\footnotesize 相邻亮纹（或暗纹）间距与波长、缝屏距、缝距关系。} & \\
  相干加强条件\newline {\footnotesize 光程差为波长的整数倍（写明暗纹则另条件）。} & \\
  相干减弱条件\newline {\footnotesize 光程差为半波长奇数倍。} & \\
  光程\newline {\footnotesize 几何路程与折射率乘积。} & \\
  洛埃德镜条纹间距\newline {\footnotesize 与双缝类似间距公式（注意半波损失对明暗互换）。} & \\
  薄膜干涉光程差\newline {\footnotesize 近似垂直入射：薄膜反射相干的两束光光程差表达式。} & \\
  半波损失\newline {\footnotesize 何时附加 $\lambda/2$ 光程（疏到密反射等表述）。} & \\
  劈尖厚度\newline {\footnotesize 小劈尖角：膜厚与水平距离、夹角关系。} & \\
  劈尖条纹间距\newline {\footnotesize 相邻条纹对应的膜厚差与条纹间距公式。} & \\
\end{tabularx}
\medskip

% --- section_15.tex ---
\begin{tabularx}{\columnwidth}{XX}
  反射定律\newline {\footnotesize 入射角等于反射角（在同一平面内）。} & \\
  折射定律\newline {\footnotesize 斯涅尔定律：入射角与折射角正弦之比。} & \\
  折射率定义\newline {\footnotesize 真空光速与介质光速之比。} & \\
  相对折射率\newline {\footnotesize 两种介质间相对折射率与各自折射率关系。} & \\
  全反射临界角\newline {\footnotesize 光密到光疏：临界角与折射率关系。} & \\
  透镜成像公式\newline {\footnotesize 薄透镜成像公式（含符号约定依教材）。} & \\
  放大率\newline {\footnotesize 像高与物高之比（或像距物距之比依符号约定）。} & \\
\end{tabularx}
\medskip

% --- section_16.tex ---
\begin{tabularx}{\columnwidth}{XX}
  热力学温度\newline {\footnotesize 摄氏温标与热力学温标换算。} & \\
  微粒数\newline {\footnotesize 物质的量与阿伏伽德罗常数、微粒数关系。} & \\
  物质的量\newline {\footnotesize 质量与摩尔质量关系。} & \\
  单个分子平均占有体积\newline {\footnotesize 固体或液体：摩尔体积与阿伏伽德罗常数。} & \\
  固体、液体分子直径估算\newline {\footnotesize 球模型：分子直径与摩尔体积关系（数量级估算）。} & \\
  气体分子平均间距估算\newline {\footnotesize 立方体模型：平均间距与摩尔体积关系。} & \\
  分子平均动能\newline {\footnotesize 温度与分子平均动能（理想气体）。} & \\
  分子平均速率\newline {\footnotesize 三种特征速率之一（写明题目常用哪一种）。} & \\
\end{tabularx}
\medskip

% --- section_17.tex ---
\begin{tabularx}{\columnwidth}{XX}
  等温变化\newline {\footnotesize 玻意耳定律：压强与体积乘积（理想气体）。} & \\
  等容变化\newline {\footnotesize 查理定律：压强与热力学温度成正比。} & \\
  等容变化比值形式\newline {\footnotesize 查理定律：初末态压强与温度比值形式。} & \\
  等压变化\newline {\footnotesize 盖—吕萨克定律：体积与热力学温度成正比。} & \\
  等压变化比值形式\newline {\footnotesize 盖—吕萨克定律：初末态体积与温度比值形式。} & \\
  理想气体状态方程\newline {\footnotesize $pV=nRT$ 或质量形式。} & \\
  液体上下压强\newline {\footnotesize 静止液体内部压强（深度、密度）。} & \\
  有质量活塞上下压强\newline {\footnotesize 活塞受力平衡：上下气体压强与活塞重力关系。} & \\
  水平液柱与活塞压强\newline {\footnotesize 水平柱塞两侧压强传递关系。} & \\
  竖直液柱加速度\newline {\footnotesize 液柱整体加速：上下压强差与液柱动力学方程。} & \\
  竖直活塞加速度\newline {\footnotesize 活塞与气体：牛顿第二定律列式后的加速度表达式。} & \\
\end{tabularx}
\medskip

% --- section_18.tex ---
\begin{tabularx}{\columnwidth}{XX}
  热力学第一定律\newline {\footnotesize $\Delta U=Q+W$（注明符号约定：外界对气体做功为正等）。} & \\
  等容过程\newline {\footnotesize 体积不变：热力学第一定律简化形式。} & \\
  等压过程气体对外做功\newline {\footnotesize 等压膨胀（或压缩）功与压强、体积变化。} & \\
  压强线性变化做功\newline {\footnotesize $p$ 随 $V$ 线性变化时气体做功（平均压强或图像面积）。} & \\
  等温过程\newline {\footnotesize 理想气体等温：内能变化为零时的吸放热与做功关系。} & \\
  绝热过程\newline {\footnotesize $Q=0$：内能变化仅由做功引起（理想气体）。} & \\
\end{tabularx}
\medskip

% --- section_19.tex ---
\begin{tabularx}{\columnwidth}{XX}
  能量守恒 & 能量既不会凭空产生，也不会凭空消失 \\
  第一类永动机 & 本质上要求无中生有地产生能量 \\
  第二类永动机 & 本质上要求热量100\% 转化为有用功 \\
\end{tabularx}
\medskip

% --- section_26.tex ---
\begin{tabularx}{\columnwidth}{XX}
  能量子关系\newline {\footnotesize 谐振子能量量子化：相邻能级间隔与频率。} & \\
  光子频率与波长\newline {\footnotesize 真空中波长与频率关系。} & \\
  光子能量表达式\newline {\footnotesize 光子能量与频率（普朗克常量）。} & \\
  振子能量量子化\newline {\footnotesize 普朗克假设：最小能量单元与频率。} & \\
  峰值波长规律\newline {\footnotesize 黑体辐射：维恩位移定律形式。} & \\
  黑体辐射规律（定性） & 温度越高，总辐射越强，峰值波长越短 \\
\end{tabularx}
\medskip

% --- section_20.tex ---
\begin{tabularx}{\columnwidth}{XX}
  光子能量\newline {\footnotesize 光子能量与频率。} & \\
  光子动量\newline {\footnotesize 光子动量与波长或频率关系。} & \\
  德布罗意波长\newline {\footnotesize 实物粒子德布罗意波长与动量。} & \\
  德布罗意关系（非相对论）\newline {\footnotesize 动量与波长的德布罗意关系。} & \\
\end{tabularx}
\medskip

% --- section_21.tex ---
\begin{tabularx}{\columnwidth}{XX}
  爱因斯坦光电方程\newline {\footnotesize 光子能量、逸出功与最大初动能关系。} & \\
  最大初动能（电压表示）\newline {\footnotesize 遏止电压与最大初动能关系。} & \\
  截止频率\newline {\footnotesize 光电效应：截止频率与逸出功。} & \\
  逸出功（波长表示）\newline {\footnotesize 用极限波长表示逸出功。} & \\
  最大初动能--频率图像\newline {\footnotesize $E_k-\nu$ 图斜率与截距的物理意义对应的公式。} & \\
  遏止电压--频率图像\newline {\footnotesize $U_c-\nu$ 图斜率与截距对应的公式。} & \\
\end{tabularx}
\medskip

% --- section_22.tex ---
\begin{tabularx}{\columnwidth}{XX}
  卢瑟福散射角判定 & 大角散射对应近核强库仑作用 \\
  玻尔能级公式（氢原子）\newline {\footnotesize 氢原子能级与量子数 $n$ 的关系。} & \\
  轨道半径（氢原子）\newline {\footnotesize 玻尔模型：第 $n$ 轨道半径与玻尔半径。} & \\
  类氢离子轨道半径\newline {\footnotesize 核电荷数为 $Z$ 时轨道半径与氢原子半径关系。} & \\
  频率条件\newline {\footnotesize 跃迁：辐射或吸收光子频率与两能级差关系。} & \\
  最大谱线数（从第 $n$ 能级向下）\newline {\footnotesize 一群氢原子从 $n$ 向下跃迁最多发射谱线条数。} & \\
\end{tabularx}
\medskip

% --- section_23.tex ---
\begin{tabularx}{\columnwidth}{XX}
  衰变规律\newline {\footnotesize 剩余核数随时间指数衰减规律。} & \\
  半衰期定义\newline {\footnotesize 半数原子核衰变所需时间（统计规律）。} & \\
  质量亏损\newline {\footnotesize 反应前后质量之差。} & \\
  质能关系\newline {\footnotesize 质量亏损与释放核能关系。} & \\
  原子质量单位换算\newline {\footnotesize $1\ \mathrm{u}$ 对应的能量（MeV）。} & \\
  结合能法放能\newline {\footnotesize 用生成物与反应物结合能之差表示释放核能。} & \\
  比结合能\newline {\footnotesize 平均每个核子的结合能定义。} & \\
\end{tabularx}
\medskip

% --- section_23.tex ---
\begin{tabularx}{\columnwidth}{|c|X|}
  \hline
  反应类型 & 反应方程式 \\
  \hline
  发现质子\newline {\footnotesize 历史上人工转变：发现质子的核反应方程。} & \\
  \hline
  发现中子\newline {\footnotesize 查德威克实验：发现中子的核反应方程。} & \\
  \hline
  发现正电子\newline {\footnotesize 人工放射性：发现正电子的核反应方程。} & \\
  \hline
  $\alpha$ 衰变\newline {\footnotesize 写出 $\alpha$ 衰变方程通式或示例。} & \\
  \hline
  $\beta^-$ 衰变本质\newline {\footnotesize 写出中子转化为质子并放出电子的本质方程。} & \\
  \hline
  核裂变\newline {\footnotesize 重核裂变典型方程（题给元素）。} & \\
  \hline
  核聚变\newline {\footnotesize 轻核聚变典型方程（如氘氚）。} & \\
  \hline
\end{tabularx}
\medskip

